\documentclass{report}
\usepackage{amsmath,amssymb}
\setlength{\parindent}{0mm}
\setlength{\parskip}{1em}
\begin{document}
\begin{center}
\rule{6in}{1pt} \
{\large Jiahao Chen \\
{\bf A Julia package for practical iterative SVDs}}

32 Vassar St \\ Rm 32-206 \\ MIT CSAIL \\ Cambridge \\ MA 02139
\\
{\tt jiahao@mit.edu}\\
Andreas Noack\\
Jake Bolewski\\
Alan Edelman\end{center}

We present an implementation of iterative Golub-Kahan-Lanczos (GKL)
bidiagonalization in the Julia package IterativeSolvers.jl, allowing for the
iterative computation of the singular value decomposition. The implementation
combines the best features of available implementations, featuring a standardized
palette of techniques aiding practical computation such as restarting and
reorthogonalization. Our implementation offers multiple restarting strategies,
such as explicit restarts, implicit restarts as in ARPACK, thick restarts as
in SLEPc, and harmonic restarts. We also offer multiple
reorthogonalization strategies, such as partial
orthogonalization as in PROPACK, and full reorthogonalization using QR and
various Gram-Schmidt methods.

We show how the multimethods (multiple dispatch) system of Julia allows
for modular composition of these
various components, allowing for very fine grained control over roundoff error.
The Julia type system also allows for bidiagonalizations to be computed over a
wide variety of matrix data types, ranging from ordinary dense matrices to
sparse matrices with multiprecision floating point numbers. Overloading elementary
operators such as multiplication (*) and subtraction (-) allow the same code to
run on new data structures, such as out of core arrays in a database, or custom
matrix representations. The same technique allows us to provide run time control
over the underlying linear algebra library, allowing users to change between
OpenBLAS, MKL, Eigen, and others. The result is a flexible yet performant
library of components which can be used for high performance computation that
at the same time facilitates further experimentation and prototyping of new
algorithms.


\end{document}
